% !PRINTED 130
\weberpage{130}
%<<CONT

setzen und erhalten, wenn wir $\alpha+h$ mit $\alpha'$ bezeichnen,
\begin{equation}
 \alpha'=\alpha-\frac{a f(\alpha)}{2Q f'(\alpha)},
\tag{11}
\end{equation}
und wenn wir $|f(\alpha')|$ mit $a'$ bezeichnen, nach (10)
\begin{equation}
 a'<a\left(1-\frac{a}{4Q}\right).
\tag{12}
\end{equation}

Wir leiten nun durch dasselbe Verfahren aus $\alpha',a'$ ein zweites Grössensystem $\alpha'',a''$ ab, dann aus $\alpha'',a''$ ein drittes $\alpha''',a'''$ u. s. f., und erhalten so für die Reihe abnehmender positiver Zahlen $a,a',a'',\ldots$ die Ungleichungen
\begin{equation}
\begin{aligned}
 a'&<a\left(1-\frac{a}{4Q}\right),\\
 a''&<a'\left(1-\frac{a'}{4Q}\right),\\
 a'''&<a''\left(1-\frac{a''}{4Q}\right),\\
 &\hspace{1.5em}\vdots
\end{aligned}
\tag{13}
\end{equation}
und beweisen nun, dass diese Zahlenreihe sich der Grenze Null nähert. Wäre dies nicht der Fall, so würde sich eine positive Zahl $\omega$ so angeben lassen, dass alle $a^{(\nu)}$ über $\omega$ bleiben; dann wären die Differenzen
\[
 \left(1-\frac{a}{4Q}\right),\quad
 \left(1-\frac{a'}{4Q}\right),\quad
 \left(1-\frac{a''}{4Q}\right)\cdots
\]
alle kleiner, als der positive echte Bruch
\[
 \Theta=1-\frac{\omega}{4Q},
\]
und aus (13) würde folgen
\[
 a'<a\Theta,\quad
 a''<a'\Theta,\quad
 a'''<a''\Theta,\quad
 \ldots,\quad
 a^{(\nu)}<a^{(\nu-1)}\Theta,
\]
woraus durch Multiplication
\begin{equation}
 a^{(\nu)}<a\Theta^\nu.
\tag{14}
\end{equation}
Darin aber liegt ein Widerspruch, da $\Theta^\nu$ mit unendlich wachsendem $\nu$ unendlich klein wird. Das so gewonnene Resultat drücken wir in der Gleichung aus
\begin{equation}
 \lim a^{(\nu)}=0.
\tag{15}
\end{equation}
